\chapter{Transferability}
\label{ch:transferability}

A framework that works only for the city it was built on is of limited scientific
value. This chapter tests research question RQ4 by transferring the entire pipeline
to two further metropolitan areas through a single parameterized code path, and
asks whether the diagnostic, generative, and validation layers all survive the move.

\section{Study design}
\label{sec:transf-design}

Two transfer cities were chosen to bracket \zmg in size and structure:
\textbf{Toluca} (Zona Metropolitana del Valle de Toluca, \num{16} municipalities,
$\sim$\num{2.3} million inhabitants, \num{538} urban \ageb{}s) as a large metro
comparable to \zmg, and \textbf{Aguascalientes} (\num{3} municipalities, \num{356}
urban \ageb{}s) as a compact one. Monterrey was the original target but its
\textsc{gtfs} feed proved unobtainable, which blocks the supply layer; it is retained
only as a Tier-1 (demand-only) reference. Both transfer cities have verified,
downloadable \textsc{gtfs}.

Every stage runs through the same code as \zmg, invoked as
\code{w9\_run\_*.py -{}-city \{tol,ags\}}, reusing the pure \zmg functions over
per-city configuration. Tier-1 inputs (census, \textsc{denue}, street network) are
national and therefore available everywhere; equity indices come from the same
CONAPO and CONEVAL products. No local Origin--Destination survey exists for either
city, so the calibrated $\beta = \num{1.2005}$ is transferred from \zmg; because the
size-comparison metrics below are ratios, they are insensitive to this choice.

\section{Demand and coverage-gap transfer}
\label{sec:transf-diagnosis}

The demand surface transfers cleanly and, more importantly, \emph{differentiates} the
three cities in interpretable ways (\cref{tab:transfer-diag}, \cref{fig:transfer}).
Mean household vehicle ownership follows a clear gradient---Toluca \num{0.529} $<$
\zmg \num{0.577} $<$ Monterrey \num{0.635} (Tier-1 reference) $<$ Aguascalientes
\num{0.667}---so the least car-dependent metro is the large, lower-income Toluca and
the most car-dependent is compact, higher-income Aguascalientes.

\begin{table}[htbp]
\centering
\caption{Diagnostic transfer across the three \textsc{gtfs} cities.}
\label{tab:transfer-diag}
\begin{tabular}{l S[table-format=4.0] S[table-format=1.3] S[table-format=2.1] S[table-format=2.1]}
\toprule
City & {Urban \ageb{}s} & {Vehicle rate} & {Unserved (\%)} & {High-gap (\%)} \\
\midrule
Guadalajara    & 1881 & 0.577 & 32.7 & 20.7 \\
Toluca         & 538  & 0.529 & 19.1 & 14.9 \\
Aguascalientes & 356  & 0.667 & 19.7 & 9.6 \\
\bottomrule
\end{tabular}
\end{table}

The coverage-gap diagnosis likewise transfers and orders the cities sensibly: the
High-gap share is highest in \zmg (\SI{20.7}{\percent}), lower in Toluca
(\SI{14.9}{\percent}), and lowest in the compact, well-covered Aguascalientes
(\SI{9.6}{\percent}). The diagnostic is thus not a \zmg artifact---it produces a
plausible, differentiated reading of unmet demand in every metro.

\begin{figure}[htbp]
  \centering
  \includegraphics[width=0.85\linewidth]{transfer_highgap.pdf}
  \caption{Diagnostic transfer: High-gap share, unserved share, and mean vehicle
    rate across the three metropolitan areas.}
  \label{fig:transfer}
\end{figure}

\section{Prioritization and corridor generation transfer}
\label{sec:transf-generation}

The Node--Place--People weighting re-derives its own weights per city rather than
importing \zmg's. Notably, \emph{social lag} (\code{pe\_rezago}) becomes the single
largest objective weight in both transfer cities (Toluca \num{0.182}, Aguascalientes
\num{0.247}), whereas in \zmg the employment and population terms led. Because the
weights are data-driven (CRITIC$\,\times\,$EWM), this shift is a feature: the method
adapts its emphasis to each city's own indicator structure.

Corridor generation produces substantive, feasible corridors in both cities
(\cref{tab:transfer-corridors}). Toluca's best feasible corridor, \code{W6\_G01}
(\SI{12.6}{\kilo\metre}, \num{23} \ageb{}s, \num{133728} trips/day, light rail), is
the direct transfer analogue of \zmg's merit-passing \code{W6\_G02}. Aguascalientes
yields \code{W6\_G05} (\SI{5.9}{\kilo\metre}, \num{12} \ageb{}s, \num{107633}
trips/day). The same feasibility-versus-directness signature seen in \zmg reappears:
in Aguascalientes a genuinely high-demand corridor (\code{W6\_G03}, \num{22}
\ageb{}s, \num{115162} trips/day) is rejected for exceeding the directness cap
(\num{2.03} > \num{1.8})---the identical behaviour documented for \zmg, now observed
in a different city.

\begin{table}[htbp]
\centering
\caption{Best feasible corridor per transfer city, with the notable infeasible case.}
\label{tab:transfer-corridors}
\begin{tabular}{ll S[table-format=2.1] S[table-format=2.0] S[table-format=6.0] c}
\toprule
City & Corridor & {km} & {\ageb{}s} & {Demand/day} & Feasible \\
\midrule
Toluca         & W6\_G01 & 12.6 & 23 & 133728 & Yes \\
Aguascalientes & W6\_G05 & 5.9  & 12 & 107633 & Yes \\
Aguascalientes & W6\_G03 & 13.9 & 22 & 115162 & {--} \\
\bottomrule
\end{tabular}
\end{table}

\section{Validation transfer}
\label{sec:transf-validation}

The audit and validation layers transfer as well, with one honest exception. The W7
audit runs on both networks (\cref{tab:transfer-audit}). Toluca's \num{622}
concessioned routes are flagged almost universally---\num{612} of \num{622}, of which
\num{431} are Redundant---quantifying the parallelism of a system with roughly thirty
operators. (Its zero feasible count is a source-\textsc{gtfs} artifact: median stop
spacing is \SI{43}{\metre}, far below the \SI{300}{\metre} constraint, so the
audit's flags rather than its feasibility gate carry the signal.) Aguascalientes,
with \num{48} routes at \SI{239}{\metre} spacing, has \num{6} feasible routes and
\num{47} flagged.

\begin{table}[htbp]
\centering
\caption{Existing-route audit (W7) for the transfer cities.}
\label{tab:transfer-audit}
\begin{tabular}{l S[table-format=3.0] S[table-format=3.0] S[table-format=3.0] S[table-format=2.0] S[table-format=2.0]}
\toprule
City & {Routes} & {Flagged} & {Redundant} & {Indirect} & {Low dem.} \\
\midrule
Toluca         & 622 & 612 & 431 & 88 & 93 \\
Aguascalientes & 48  & 47  & 4   & 33 & 10 \\
\bottomrule
\end{tabular}
\end{table}

The W8 benchmark reveals a genuine transferability insight. In both transfer cities
the generated corridors overlap the existing network \emph{heavily}---Toluca
\SI{75.0}{\percent}, Aguascalientes \SI{54.3}{\percent}---the opposite of \zmg's low
overlap. In these already well-served networks ($\sim$\SI{80}{\percent} baseline
coverage) the highest-demand corridors are already served, so the generator
\emph{re-identifies} them (revealed-preference corroboration) rather than opening new
coverage; correspondingly, the before/after gains are modest (Toluca coverage
$+\num{0.4}$~pp, Gini $-\num{0.0031}$; Aguascalientes $+\num{0.3}$~pp,
$-\num{0.0027}$).

The mask-and-reconstruct backtest, however, does \emph{not} transfer to a bus-only
network. Neither city has a premium tier to hold out. A demand-trunk proxy applied to
Toluca (masking the \num{23} busiest routes) collapses: frontier anchors fall from
\num{14} to \num{6} and no multi-anchor corridor can form, so nothing is re-proposed.
The mechanism is intrinsic---\zmg's premium rail is a separable overlay redundant with
parallel buses, so masking it leaves the served/unserved seam intact, whereas a bus
trunk \emph{is} the local service and masking it erases the seam the generator depends
on. This is a documented precondition of the backtest, not a defect of the framework;
the benchmark and before/after metrics remain the operative validation for cities
without a premium tier.

\section{Transfer error and limitations}
\label{sec:transf-limits}

Four caveats bound the transfer claim. First, \textbf{no local calibration}: $\beta$
is transferred from \zmg, so absolute demand magnitudes carry the prior's error,
though the ratio-based diagnostics do not. Second, \textbf{a coarser equity input}:
CONEVAL publishes only a categorical social-lag \emph{grade} at \ageb level for 2020,
so \code{pe\_rezago} is an ordinal approximation of \zmg's continuous index
(direction preserved, granularity coarser). Third, \textbf{unit-scale incomparability}:
the \ageb is much coarser relative to population in semi-rural Toluca, so cross-city
\ageb counts are not directly comparable and only per-\ageb and per-capita metrics
are. Fourth, as shown above, \textbf{the hold-out backtest requires a premium tier}
that bus-only networks lack. None of these undermines the central result: the
complete framework transfers end-to-end to two metros of very different size through
one code path, on cities Monterrey's missing feed had made impossible.
